% =====================================================================
% 04 — Which objections actually hold?
% =====================================================================
% Planning (% comments only):
% For each claim from sec 3, deliver a verdict in {theorem / robust
% empirical / weak empirical / refuted / untested}. Each subsection
% has the same structure: claim, premises, what supports it, what
% would refute it, current verdict, what the companion suite measures.
% This is the technical core. Be specific.
% =====================================================================

\section{Which Objections Actually Hold}
\label{sec:what_holds}

We now grade each claim from \cref{sec:catalog} against a clear bar:
a claim is a \emph{theorem} only if it follows from the math of the factorization, parameterization, or objective;
a \emph{robust empirical} finding only if it is reproduced across labs and scales with a stated mechanism;
\emph{weak empirical} if its support depends on a single setup or a fragile premise;
and \emph{refuted} if a direct counterexample exists under conditions a deployed system can satisfy.
\Cref{tab:verdicts} previews the verdicts;
the subsections justify each.

\begin{table}[ht]
  \centering
  \small
  \caption{Standing of each objection.
  ``Layer'' is from \cref{sec:catalog}.
  ``Verdict'' is the standing of the claim under the bar of theorem / robust empirical / weak empirical / refuted / untested.
  ``Bites'' is what the surviving form of the claim indicts.}
  \label{tab:verdicts}
  \begin{tabular}{@{}lllll@{}}
    \toprule
    Tag & Claim & Layer & Verdict & Bites \\
    \midrule
    \ref{cl:softmax_bottleneck}  & Softmax rank cap                     & architecture     & theorem            & softmax head \\
    \ref{cl:rank_collapse}       & Pure-attention rank collapse         & architecture     & theorem (pure); masked (trained) & softmax attention \\
    \ref{cl:partition_local}     & Local partition function tax         & architecture     & established        & softmax / normalization \\
    \ref{cl:fixed_compute}       & Fixed compute ceiling                & architecture     & theorem            & one-pass generation \\
    \ref{cl:hidden_state}        & Fixed-state recall limit             & architecture     & theorem (SSMs); workaround (attn) & fixed-state heads \\
    \ref{cl:mode_covering}       & Mode-covering MLE                    & objective        & theorem            & MLE \\
    \ref{cl:teacher_forcing}     & Exposure bias                        & objective        & weak empirical     & objective (small effect at scale) \\
    \ref{cl:next_token_misalign} & Likelihood $\neq$ correctness        & objective        & true and central   & objective \\
    \ref{cl:compounding}         & $(1-\varepsilon)^T$ compounding      & trained behavior & valid under premises; refuted with a verifier & blind rollout only \\
    \ref{cl:reversal}            & Reversal curse                       & trained behavior & robust empirical   & directional training \\
    \ref{cl:hallucination}       & Hallucination is inevitable          & trained behavior & weak (strong form refuted by abstention + retrieval) & MLE + open generation \\
    \ref{cl:lipschitz_margin}    & Lipschitz--margin brittleness        & trained behavior & robust empirical + theory link & unconstrained Lipschitz \\
    \ref{cl:data_wall}           & Data wall                            & external         & robust empirical (resource curve) & data-bound scaling \\
    \bottomrule
  \end{tabular}
\end{table}

\subsection{Theorems: rank, collapse, partition, compute, mode-covering}
\label{ssec:theorems}

The strongest objections to the standard parameterization and objective are not conjectures.
They follow from the math.

\paragraph{\ref{cl:softmax_bottleneck} Softmax bottleneck.}
\citet{yang2018softmax} show that a softmax over a $d$-dimensional hidden state can only produce a stochastic matrix of context-conditional probabilities whose log-probability matrix has rank at most $d+1$.
If natural language has higher rank in this sense --- as their empirical analysis suggests --- the model class is misspecified by construction.
Mixture-of-softmaxes restores rank;
\citet{kanai2018mixture} and \citet{chang2022softmax} confirm the predicted improvement and identify mechanisms behind it.
This is a theorem.
It indicts the softmax parameterization, not the autoregressive factorization.

\paragraph{\ref{cl:rank_collapse} Rank collapse.}
\citet{dong2021attention} prove that without residuals, MLPs, or layer normalization, a stack of self-attention layers contracts representations to a rank-one matrix at a doubly-exponential rate in depth.
The proof identifies the row-stochastic attention matrix as a Perron--Frobenius contraction on the simplex.
The theorem is exact for pure attention.
In trained transformers the effect is observed but masked by skip connections and feedforward blocks~\citep{noci2022signal,he2023deep};
whether residual streams alone are enough remains an open empirical question.
For the present paper, the cleanest statement is: pure attention provably collapses;
the contribution of residual stream depth to recovery is quantitative and architecture-specific.

\paragraph{\ref{cl:partition_local} Local partition function tax.}
Softmax attention is, by \citet{ramsauer2020hopfield}, a single gradient step on a continuous Hopfield energy whose exponential form was chosen because it admits a closed-form normalizer.
Most downstream uses of the attention layer need only \emph{differences} between scores, in which case the local $Z$ cancels in expectation.
Whether the per-layer normalization is necessary --- as opposed to being a convenient bookkeeping device --- is itself a falsifiable architectural claim, and is exactly what sigmoid attention~\citep{wortsman2023sigmoid,ramapuram2024theory} and energy heads~\citep{grathwohl2019your} probe.
We mark this claim as ``established'' rather than ``theorem'' because what is established is the cost, not the necessity.

\paragraph{\ref{cl:fixed_compute} Fixed compute ceiling.}
A standard transformer with depth $L$, width $w$, and head count $h$ performs a bounded number of sequential operations per generated token.
\citet{merrill2023parallelism,strobl2024transformer} formalize this within circuit complexity: transformers without intermediate steps are placed in fixed-depth uniform classes that cannot solve problems requiring more sequential depth.
Chain-of-thought~\citep{wei2022chain} and \citet{feng2023towards}'s analysis of intermediate steps explicitly use the output token stream to buy more sequential computation.
The bound is real.
The empirical question is how much chain-of-thought, search, and tool use already buy back, and at what cost in latency.

\paragraph{\ref{cl:mode_covering} Mode-covering MLE.}
For data distribution $\poftxt$ and model $\ptheta$, the forward KL is
\begin{equation}
  \KL(\poftxt \,\|\, \ptheta)
  \;=\; \EE_{\poftxt}\!\left[\log \poftxt(x) - \log \ptheta(x)\right],
\end{equation}
which diverges whenever $\ptheta(x) \to 0$ while $\poftxt(x) > 0$.
The optimizer is therefore compelled to spread mass across every observed mode.
\citet{theis2016note} and \citet{minka2005divergence} state this as the standard consequence of forward versus reverse KL;
under the noise present in any web-scale corpus, the model assigns nonzero mass to contradictions and incorrect continuations by construction.
This is a theorem about the objective.
It indicts MLE and, indirectly, the family of methods that use MLE as a pretraining proxy without correction --- though instruction tuning and reinforcement learning from verifier feedback are exactly the correction stage.

\subsection{The $(1-\varepsilon)^T$ argument, carefully}
\label{ssec:compounding}

The $(1-\varepsilon)^T$ argument is the most-cited critique of autoregressive language models in popular discussion.
We dwell on it because the gap between its mathematical validity and its actual reach is the main case study for layer discipline in this paper.

The argument is exact under three explicit premises.
Let $X_t$ denote the model's $t$-th generated token and $X_t^\star$ denote the (multi-set of) correct alternatives.
\begin{itemize}
  \item \emph{Premise 1 (constant rate).}
        There exists $\varepsilon \in (0,1)$ such that $\PP(X_t \neq X_t^\star) = \varepsilon$ for every $t$.
  \item \emph{Premise 2 (independence).}
        The events $\{X_t = X_t^\star\}_{t=1}^{T}$ are mutually independent.
  \item \emph{Premise 3 (unrecoverability).}
        Once $X_t \neq X_t^\star$, the trajectory cannot recover; equivalently, ``correct'' means literal token-wise match against a single reference path.
\end{itemize}
Under premises 1--3, $\PP[\text{trajectory correct}] = \prod_t (1-\varepsilon) = (1-\varepsilon)^T$.
For any $\varepsilon > 0$, this decays exponentially in $T$~\citep{lecun2022autonomous,bachmann2024pitfalls}.

The math is correct.
The three premises do the work.

Premise 1 fails in trained models: the realized per-token error rate depends on context, position, topic, and prompt;
some stretches are nearly deterministic and others are unstable~\citep{kalai2024calibrated}.
A constant-rate model is an envelope calculation, not a description of the empirical process.

Premise 2 also fails.
Errors cluster.
A wrong token typically conditions a region of corrupted state in which subsequent tokens are more, not equally, likely to be wrong.
For an actual i.i.d. claim one would have to control for context, and the literature offering such controls reports a context-dependent error rate rather than independent errors~\citep{stechly2024chain,lin2024diverging}.

Premise 3 is the most consequential.
A wrong token is unrecoverable only if the surrounding system has no mechanism to detect, reject, or revise it.
Modern systems are designed precisely around the opposite assumption.
A compiler rejects a syntactically invalid program;
a Lean kernel rejects an invalid proof step~\citep{moura2021lean};
a unit-test harness flags a behavioral failure;
a search procedure can branch and backtrack;
a verifier-in-the-loop trainer reinforces only the trajectories that survive the checker~\citep{lightman2024lets,uesato2022solving}.
A formal-methods workflow can sit in a generate--check--revise loop that turns a long trajectory into a sequence of short, verifier-bounded segments.

The interesting empirical object is not $\PP[\text{token locally wrong}]$.
It is closer to
\begin{equation}
  \PP[\text{a consequential error survives detection, recovery, and verification}].
  \label{eq:effective_eps}
\end{equation}
That quantity can be much smaller than $\varepsilon$, especially in domains with a hard checker.
It can also still be too large.
The point is that \cref{eq:effective_eps} is a different object from $(1-\varepsilon)^T$.
The compounding bound applies to blind rollout.
It does not apply to systems whose construction is the negation of premises 1--3.

The companion suite (\texttt{experiments/02\_ar\_pros\_cons/probes/probe\_06\_error\_compounding.py}) pre-registers a recoverable-Markov alternative model whose null is geometric decay and whose alternative is a slower-than-geometric process with verifier-bounded segments.
Whether the realized rate in formal-methods generation is geometric is the empirical question we explicitly defer to that suite.

\subsection{Robust empirical and weak empirical findings}
\label{ssec:empirical}

\paragraph{\ref{cl:reversal} Reversal curse (robust empirical).}
\citet{berglund2023reversal} demonstrate that models trained on ``$A$ is $B$'' do not reliably retrieve ``$B$ is $A$''.
The result reproduces across model scales and across structured-relation tests.
The indicted target is the directional next-token objective combined with static text training, not the autoregressive factorization itself --- as the experimental controls show, conditioning the same model on a paraphrase that recovers symmetric structure can lift performance.

\paragraph{\ref{cl:lipschitz_margin} Lipschitz--margin (robust empirical + theory link).}
\citet{tsuzuku2018lipschitz} and \citet{virmaux2018lipschitz} give the explicit bound and the upper estimate $L \leq \prod_i \|W_i\|_2$.
Adversarial robustness studies~\citep{szegedy2014intriguing,goodfellow2015explaining,carlini2023llm} report small input perturbations that swing model outputs across decision boundaries despite high confidence.
The mechanism is consistent: confidence and geometric margin are decoupled.
The indictment is unconstrained Lipschitz training rather than autoregression.

\paragraph{\ref{cl:hallucination} Hallucination (weak in its strong form).}
\citet{kalai2024calibrated} prove that calibrated next-token training places nonzero mass on incorrect continuations and conclude that an open-generation system trained this way will hallucinate at a nonzero rate.
The strong version of the claim --- hallucination is \emph{inevitable} --- requires also that the system not abstain, not retrieve, and not check.
\citet{kang2023impact} and \citet{maynez2020faithfulness} document realized hallucination patterns;
\citet{liu2024lost} document grounding effects.
We take the conservative reading: the structural pressure follows from the objective, but the realized rate is a system property that retrieval, abstention, and verification can change~\citep{ji2023survey}.

\paragraph{\ref{cl:teacher_forcing} Exposure bias (weak empirical).}
Scheduled sampling~\citep{bengio2015scheduled} and reinforcement-learning extensions~\citep{ranzato2016sequence} demonstrated the mechanism on small models.
On large models, the realized gap between teacher-forced training and free-running generation appears smaller~\citep{wang2020exposure}, though the question is not closed for very long horizons.
We flag this as weak rather than refuted because its magnitude in the long-horizon regime is precisely the regime where it would matter most, and the existing evidence is concentrated in short-horizon benchmarks.

\paragraph{\ref{cl:data_wall} Data wall (robust empirical, resource curve).}
The argument depends only on (i) the empirical scaling law that loss requires exponential data for linear improvement and (ii) an estimate of available high-quality text~\citep{villalobos2024position}.
Neither piece requires a ``model cannot do X'' claim;
both are observed.
The companion suite includes a probe that re-fits the scaling law on \textsc{VeriBench} subsets~\citep{daruru2024veribench} and extrapolates.

\subsection{Explicitly excluded as not load-bearing}
\label{ssec:excluded}

Two further objections appear regularly and are not load-bearing in our reading.

The ``stochastic parrot'' framing~\citep{bender2021dangers} is not operationalizable in its general form: it makes no falsifiable prediction at the granularity at which it is stated.
Its operationalizable pieces (reversal, planning, compositional generalization) reduce to claims we have already addressed.

The claim that autoregressive factorization \emph{itself} caps expressivity is mostly false.
The chain rule of probability decomposes any joint over a finite sequence into left-to-right conditionals.
Practical bottlenecks come from the parameterization of the conditional, the objective used to train it, and the compute used to run inference --- not from the existence of the factorization.

\subsection{What this leaves}
\label{ssec:what_remains}

The surviving load-bearing claims separate cleanly into three families:
\begin{enumerate}
  \item \emph{Against the softmax / local normalization:} \ref{cl:softmax_bottleneck}, \ref{cl:rank_collapse}, \ref{cl:partition_local}.
  \item \emph{Against the MLE objective:} \ref{cl:mode_covering}.
  \item \emph{Against fixed-compute generation:} \ref{cl:fixed_compute}, \ref{cl:lipschitz_margin}.
\end{enumerate}
The next three sections take this as input and ask, for each family, what alternative architectures put in its place.
The recurring question is whether the alternative pays the same cost in a different place or absorbs it.
